% --- Archon physics formalization source begin ---
% archon:physics
% archon:covers PhyXMiniProblems/problem_phyx_mini_0596.lean
% archon:source-report reports/phyx_mini/problem_phyx_mini_0596.source.json

\chapter{Physics problem phyx\_mini\_0596}
\label{ch:PhyXMiniProblems_problem_phyx_mini_0596}

\paragraph{Problem source.}
\#\# Physical scenario

Figure is a graph of intensity versus wavelength for light reaching Earth from galaxy NGC 7319, which is about \textbackslash{}( 3 \textbackslash{}times 10\textasciicircum{}8 \textbackslash{}) light-years away. The most intense light is emitted by the oxygen in NGC 7319. In a laboratory that emission is at wavelength \textbackslash{}( \textbackslash{}lambda = 513\textbackslash{},\textbackslash{}text\{nm\} \textbackslash{}), but in the light from NGC 7319 it has been shifted to \textbackslash{}( 525\textbackslash{},\textbackslash{}text\{nm\} \textbackslash{}) due to the Doppler effect (all the emissions from NGC 7319 have been shifted).

\#\# Question

What is the radial speed of NGC 7319 relative to Earth?

\#\# Answer choices

- A: A: 3.00 \textbackslash{}times 10\textasciicircum{}5\textbackslash{}

- B: B: 7.00 \textbackslash{}times 10\textasciicircum{}7\textbackslash{}

- C: C: 7.00 \textbackslash{}times 10\textasciicircum{}6\textbackslash{}

- D: D: 2.00 \textbackslash{}times 10\textasciicircum{}6\textbackslash{}

\#\# Figure caption (auxiliary; use the image as primary evidence)

The image is a graph showing the intensity of light as a function of wavelength, measured in nanometers (nm). 

- **Axes**: 
  - The x-axis is labeled "Wavelength (nm)" and ranges from 400 to 750 nm.
  - The y-axis is labeled "Intensity" with values from 0 to 800, though there's a break in the axis at lower intensities.

- **Data**: 
  - A green line represents the intensity data across the given wavelength range. It shows notable peaks and fluctuations.
  - Major peaks appear around 500 nm and 650 nm.

- **Annotations**:
  - A label pointing to the left of the 500 nm peak reads "Laboratory wavelength."
  - An arrow indicates a shift to the right of the 500 nm peak with a label \textbackslash{}(\textbackslash{}Delta \textbackslash{}lambda = +12 \textbackslash{}, \textbackslash{}text\{nm\}\textbackslash{}).
  - Another label identifies the graph as representing "NGC 7319."

- **Grid**: The graph has a grid to aid in reading the data. 

This graph appears to depict a spectrum, possibly displaying redshift or blueshift phenomena relevant to astrophysics.

\#\# Dataset metadata

Relativity; Physical Model Grounding Reasoning, Multi-Formula Reasoning

\paragraph{Recorded answer/context.}
C: C: 7.00 \textbackslash{}times 10\textasciicircum{}6\textbackslash{}

\paragraph{Figure/image path.}
/root/proposal\_for\_physic/science-mango/phyx\_mini\_run/phyx\_data/test\_image/596.png

\paragraph{Formalization target.}
create a compiling Lean file with sorry bodies at `PhyXMiniProblems/problem\_phyx\_mini\_0596.lean`. The Lean declarations must preserve the physical quantities, dimensions or dimensional roles, figure labels, governing-law hypotheses, and final relation expressed by this problem.
Use Mathlib/Physlib names found through LeanExplore where available. If a physics API is missing, introduce faithful local abstractions rather than scalar placeholder aliases.

\begin{theorem}[Physics formalization target]
\label{thm:physics:phyx_mini_0596:target}
\lean{PhyXMiniProblems.ProblemPhyXMini0596.radialSpeedOfNGC7319_is_choiceC}
\uses{def:physics:phyx-mini-0596:phyxminiproblems-problemphyxmini0596-speedinmeterspersecond, def:physics:phyx-mini-0596:phyxminiproblems-problemphyxmini0596-vacuumspeedoflightinmeterspersecond, def:physics:phyx-mini-0596:phyxminiproblems-problemphyxmini0596-galaxyredshiftsetup, def:physics:phyx-mini-0596:phyxminiproblems-problemphyxmini0596-matchesngc7319spectrumscenario, def:physics:phyx-mini-0596:phyxminiproblems-problemphyxmini0596-matchesprimaryspectrumfigure, def:physics:phyx-mini-0596:phyxminiproblems-problemphyxmini0596-matchesngc7319problemreadouts, def:physics:phyx-mini-0596:phyxminiproblems-problemphyxmini0596-hasphysicalredshiftparameters, def:physics:phyx-mini-0596:phyxminiproblems-problemphyxmini0596-satisfieslowspeedradialopticaldopplerlaw, def:physics:phyx-mini-0596:phyxminiproblems-problemphyxmini0596-isuniqueclosestradialspeedchoice}
The assigned autoformalize agent should translate this physics problem into Lean declarations in the covered file, with theorem and lemma proof bodies written as `by sorry`.
\end{theorem}
\begin{proof}
This is an autoformalization task, not a proof task. Produce statements that can later be proved by the physics prover without weakening the physical model.
\end{proof}

% --- Archon declaration topology begin ---
\section{Declaration topology}
\begin{definition}[Length Quantity]
  \label{def:physics:phyx-mini-0596:phyxminiproblems-problemphyxmini0596-lengthquantity}
  \lean{PhyXMiniProblems.ProblemPhyXMini0596.LengthQuantity}
  A nonnegative, unit-independent physical length
\end{definition}
\begin{proof}
  This is the declared model definition of Length Quantity.
\end{proof}
\begin{definition}[Speed Quantity]
  \label{def:physics:phyx-mini-0596:phyxminiproblems-problemphyxmini0596-speedquantity}
  \lean{PhyXMiniProblems.ProblemPhyXMini0596.SpeedQuantity}
  A nonnegative, unit-independent physical speed
\end{definition}
\begin{proof}
  This is the declared model definition of Speed Quantity.
\end{proof}
\begin{definition}[length Readout]
  \label{def:physics:phyx-mini-0596:phyxminiproblems-problemphyxmini0596-lengthreadout}
  \lean{PhyXMiniProblems.ProblemPhyXMini0596.lengthReadout}
  \uses{def:physics:phyx-mini-0596:phyxminiproblems-problemphyxmini0596-lengthquantity}
  Read a physical length in a selected length unit
\end{definition}
\begin{proof}
  This is the declared model definition of length Readout.
\end{proof}
\begin{definition}[wavelength In Nanometers]
  \label{def:physics:phyx-mini-0596:phyxminiproblems-problemphyxmini0596-wavelengthinnanometers}
  \lean{PhyXMiniProblems.ProblemPhyXMini0596.wavelengthInNanometers}
  \uses{def:physics:phyx-mini-0596:phyxminiproblems-problemphyxmini0596-lengthquantity, def:physics:phyx-mini-0596:phyxminiproblems-problemphyxmini0596-lengthreadout}
  Nanometre readout of a wavelength
\end{definition}
\begin{proof}
  This is the declared model definition of wavelength In Nanometers.
\end{proof}
\begin{definition}[distance In Light Years]
  \label{def:physics:phyx-mini-0596:phyxminiproblems-problemphyxmini0596-distanceinlightyears}
  \lean{PhyXMiniProblems.ProblemPhyXMini0596.distanceInLightYears}
  \uses{def:physics:phyx-mini-0596:phyxminiproblems-problemphyxmini0596-lengthquantity, def:physics:phyx-mini-0596:phyxminiproblems-problemphyxmini0596-lengthreadout}
  Light-year readout of an astronomical distance
\end{definition}
\begin{proof}
  This is the declared model definition of distance In Light Years.
\end{proof}
\begin{definition}[speed In Meters Per Second]
  \label{def:physics:phyx-mini-0596:phyxminiproblems-problemphyxmini0596-speedinmeterspersecond}
  \lean{PhyXMiniProblems.ProblemPhyXMini0596.speedInMetersPerSecond}
  \uses{def:physics:phyx-mini-0596:phyxminiproblems-problemphyxmini0596-speedquantity}
  Metres-per-second readout of a nonnegative physical speed
\end{definition}
\begin{proof}
  This is the declared model definition of speed In Meters Per Second.
\end{proof}
\begin{definition}[vacuum Speed Of Light In Meters Per Second]
  \label{def:physics:phyx-mini-0596:phyxminiproblems-problemphyxmini0596-vacuumspeedoflightinmeterspersecond}
  \lean{PhyXMiniProblems.ProblemPhyXMini0596.vacuumSpeedOfLightInMetersPerSecond}
  Physlib's exact vacuum light speed, read in metres per second
\end{definition}
\begin{proof}
  This is the declared model definition of vacuum Speed Of Light In Meters Per Second.
\end{proof}
\begin{definition}[Galaxy Label]
  \label{def:physics:phyx-mini-0596:phyxminiproblems-problemphyxmini0596-galaxylabel}
  \lean{PhyXMiniProblems.ProblemPhyXMini0596.GalaxyLabel}
  The emitting galaxy named in the problem and image
\end{definition}
\begin{proof}
  This is the declared model definition of Galaxy Label.
\end{proof}
\begin{definition}[Reference Frame]
  \label{def:physics:phyx-mini-0596:phyxminiproblems-problemphyxmini0596-referenceframe}
  \lean{PhyXMiniProblems.ProblemPhyXMini0596.ReferenceFrame}
  The two inertial-frame roles needed for the radial measurement
\end{definition}
\begin{proof}
  This is the declared model definition of Reference Frame.
\end{proof}
\begin{definition}[Radial Direction]
  \label{def:physics:phyx-mini-0596:phyxminiproblems-problemphyxmini0596-radialdirection}
  \lean{PhyXMiniProblems.ProblemPhyXMini0596.RadialDirection}
  Whether the source is approaching or receding along the line of sight
\end{definition}
\begin{proof}
  This is the declared model definition of Radial Direction.
\end{proof}
\begin{definition}[Emitting Species]
  \label{def:physics:phyx-mini-0596:phyxminiproblems-problemphyxmini0596-emittingspecies}
  \lean{PhyXMiniProblems.ProblemPhyXMini0596.EmittingSpecies}
  Chemical species explicitly identified as producing the strongest line
\end{definition}
\begin{proof}
  This is the declared model definition of Emitting Species.
\end{proof}
\begin{definition}[Spectrum Axis Label]
  \label{def:physics:phyx-mini-0596:phyxminiproblems-problemphyxmini0596-spectrumaxislabel}
  \lean{PhyXMiniProblems.ProblemPhyXMini0596.SpectrumAxisLabel}
  The two labeled coordinate axes of the supplied spectrum
\end{definition}
\begin{proof}
  This is the declared model definition of Spectrum Axis Label.
\end{proof}
\begin{definition}[Spectrum Trace Color]
  \label{def:physics:phyx-mini-0596:phyxminiproblems-problemphyxmini0596-spectrumtracecolor}
  \lean{PhyXMiniProblems.ProblemPhyXMini0596.SpectrumTraceColor}
  The green trace color visible in the primary image
\end{definition}
\begin{proof}
  This is the declared model definition of Spectrum Trace Color.
\end{proof}
\begin{definition}[Galaxy Spectrum Figure]
  \label{def:physics:phyx-mini-0596:phyxminiproblems-problemphyxmini0596-galaxyspectrumfigure}
  \lean{PhyXMiniProblems.ProblemPhyXMini0596.GalaxySpectrumFigure}
  \uses{def:physics:phyx-mini-0596:phyxminiproblems-problemphyxmini0596-galaxylabel, def:physics:phyx-mini-0596:phyxminiproblems-problemphyxmini0596-spectrumaxislabel, def:physics:phyx-mini-0596:phyxminiproblems-problemphyxmini0596-spectrumtracecolor}
  Literal and qualitative evidence from phyx\_data/test\_image/596.png. Intensity numbers are scalar graph readouts because the vertical axis supplies no physical unit. This structure contains no radial-speed value or answer choice
\end{definition}
\begin{proof}
  This is the declared model definition of Galaxy Spectrum Figure.
\end{proof}
\begin{definition}[Galaxy Redshift Setup]
  \label{def:physics:phyx-mini-0596:phyxminiproblems-problemphyxmini0596-galaxyredshiftsetup}
  \lean{PhyXMiniProblems.ProblemPhyXMini0596.GalaxyRedshiftSetup}
  \uses{def:physics:phyx-mini-0596:phyxminiproblems-problemphyxmini0596-lengthquantity, def:physics:phyx-mini-0596:phyxminiproblems-problemphyxmini0596-speedquantity, def:physics:phyx-mini-0596:phyxminiproblems-problemphyxmini0596-galaxylabel, def:physics:phyx-mini-0596:phyxminiproblems-problemphyxmini0596-referenceframe, def:physics:phyx-mini-0596:phyxminiproblems-problemphyxmini0596-radialdirection, def:physics:phyx-mini-0596:phyxminiproblems-problemphyxmini0596-emittingspecies, def:physics:phyx-mini-0596:phyxminiproblems-problemphyxmini0596-galaxyspectrumfigure}
  The physical spectrum and the requested observable. SpectralLine is kept abstract so that the statement that all emissions are redshifted can be expressed without pretending that a spectrum has only the oxygen peak. The plotted intensity is explicitly a scalar readout rather than a scalar replacement for a physical intensity quantity
\end{definition}
\begin{proof}
  This is the declared model definition of Galaxy Redshift Setup.
\end{proof}
\begin{definition}[Matches NGC7319 Spectrum Scenario]
  \label{def:physics:phyx-mini-0596:phyxminiproblems-problemphyxmini0596-matchesngc7319spectrumscenario}
  \lean{PhyXMiniProblems.ProblemPhyXMini0596.MatchesNGC7319SpectrumScenario}
  \uses{def:physics:phyx-mini-0596:phyxminiproblems-problemphyxmini0596-wavelengthinnanometers, def:physics:phyx-mini-0596:phyxminiproblems-problemphyxmini0596-galaxyredshiftsetup}
  Object, frame, emitting-species, dominance, and common-redshift information stated in the prose. The unknown speed is given only its physical role and direction, not a numerical value
\end{definition}
\begin{proof}
  This is the declared model definition of Matches NGC7319 Spectrum Scenario.
\end{proof}
\begin{definition}[Matches Primary Spectrum Figure]
  \label{def:physics:phyx-mini-0596:phyxminiproblems-problemphyxmini0596-matchesprimaryspectrumfigure}
  \lean{PhyXMiniProblems.ProblemPhyXMini0596.MatchesPrimarySpectrumFigure}
  \uses{def:physics:phyx-mini-0596:phyxminiproblems-problemphyxmini0596-galaxyredshiftsetup}
  Salient labels, bounds, color, and annotation visible in the raster
\end{definition}
\begin{proof}
  This is the declared model definition of Matches Primary Spectrum Figure.
\end{proof}
\begin{definition}[Matches NGC7319 Problem Readouts]
  \label{def:physics:phyx-mini-0596:phyxminiproblems-problemphyxmini0596-matchesngc7319problemreadouts}
  \lean{PhyXMiniProblems.ProblemPhyXMini0596.MatchesNGC7319ProblemReadouts}
  \uses{def:physics:phyx-mini-0596:phyxminiproblems-problemphyxmini0596-wavelengthinnanometers, def:physics:phyx-mini-0596:phyxminiproblems-problemphyxmini0596-distanceinlightyears, def:physics:phyx-mini-0596:phyxminiproblems-problemphyxmini0596-galaxyredshiftsetup}
  The numerical data supplied in the problem and its figure: distance 3 × 10\textasciicircum{}8 ly, laboratory oxygen wavelength 513 nm, observed oxygen wavelength 525 nm, and their +12 nm difference. No speed occurs here
\end{definition}
\begin{proof}
  This is the declared model definition of Matches NGC7319 Problem Readouts.
\end{proof}
\begin{definition}[Has Physical Redshift Parameters]
  \label{def:physics:phyx-mini-0596:phyxminiproblems-problemphyxmini0596-hasphysicalredshiftparameters}
  \lean{PhyXMiniProblems.ProblemPhyXMini0596.HasPhysicalRedshiftParameters}
  \uses{def:physics:phyx-mini-0596:phyxminiproblems-problemphyxmini0596-wavelengthinnanometers, def:physics:phyx-mini-0596:phyxminiproblems-problemphyxmini0596-distanceinlightyears, def:physics:phyx-mini-0596:phyxminiproblems-problemphyxmini0596-speedinmeterspersecond, def:physics:phyx-mini-0596:phyxminiproblems-problemphyxmini0596-vacuumspeedoflightinmeterspersecond, def:physics:phyx-mini-0596:phyxminiproblems-problemphyxmini0596-galaxyredshiftsetup}
  Positivity conditions needed for wavelength ratios and the physical speed
\end{definition}
\begin{proof}
  This is the declared model definition of Has Physical Redshift Parameters.
\end{proof}
\begin{definition}[Satisfies Low Speed Radial Optical Doppler Law]
  \label{def:physics:phyx-mini-0596:phyxminiproblems-problemphyxmini0596-satisfieslowspeedradialopticaldopplerlaw}
  \lean{PhyXMiniProblems.ProblemPhyXMini0596.SatisfiesLowSpeedRadialOpticalDopplerLaw}
  \uses{def:physics:phyx-mini-0596:phyxminiproblems-problemphyxmini0596-wavelengthinnanometers, def:physics:phyx-mini-0596:phyxminiproblems-problemphyxmini0596-speedinmeterspersecond, def:physics:phyx-mini-0596:phyxminiproblems-problemphyxmini0596-vacuumspeedoflightinmeterspersecond, def:physics:phyx-mini-0596:phyxminiproblems-problemphyxmini0596-galaxyredshiftsetup}
  The standard low-speed radial optical Doppler model, applied uniformly to all emission lines from one source. It relates the independently stored unknown radial speed to wavelength observations but does not state its calculated value or select an answer choice
\end{definition}
\begin{proof}
  This is the declared model definition of Satisfies Low Speed Radial Optical Doppler Law.
\end{proof}
\begin{definition}[Answer Choice]
  \label{def:physics:phyx-mini-0596:phyxminiproblems-problemphyxmini0596-answerchoice}
  \lean{PhyXMiniProblems.ProblemPhyXMini0596.AnswerChoice}
  Labels of the four radial-speed choices printed with the problem
\end{definition}
\begin{proof}
  This is the declared model definition of Answer Choice.
\end{proof}
\begin{definition}[meters Per Second]
  \label{def:physics:phyx-mini-0596:phyxminiproblems-problemphyxmini0596-answerchoice-meterspersecond}
  \lean{PhyXMiniProblems.ProblemPhyXMini0596.AnswerChoice.metersPerSecond}
  \uses{def:physics:phyx-mini-0596:phyxminiproblems-problemphyxmini0596-answerchoice}
  Displayed candidate radial speeds, in metres per second
\end{definition}
\begin{proof}
  This is the declared model definition of meters Per Second.
\end{proof}
\begin{definition}[answer Choice Error Meters Per Second]
  \label{def:physics:phyx-mini-0596:phyxminiproblems-problemphyxmini0596-answerchoiceerrormeterspersecond}
  \lean{PhyXMiniProblems.ProblemPhyXMini0596.answerChoiceErrorMetersPerSecond}
  \uses{def:physics:phyx-mini-0596:phyxminiproblems-problemphyxmini0596-speedinmeterspersecond, def:physics:phyx-mini-0596:phyxminiproblems-problemphyxmini0596-galaxyredshiftsetup, def:physics:phyx-mini-0596:phyxminiproblems-problemphyxmini0596-answerchoice}
  Absolute error between the physical radial speed and a displayed choice
\end{definition}
\begin{proof}
  This is the declared model definition of answer Choice Error Meters Per Second.
\end{proof}
\begin{definition}[Is Unique Closest Radial Speed Choice]
  \label{def:physics:phyx-mini-0596:phyxminiproblems-problemphyxmini0596-isuniqueclosestradialspeedchoice}
  \lean{PhyXMiniProblems.ProblemPhyXMini0596.IsUniqueClosestRadialSpeedChoice}
  \uses{def:physics:phyx-mini-0596:phyxminiproblems-problemphyxmini0596-galaxyredshiftsetup, def:physics:phyx-mini-0596:phyxminiproblems-problemphyxmini0596-answerchoice, def:physics:phyx-mini-0596:phyxminiproblems-problemphyxmini0596-answerchoiceerrormeterspersecond}
  A displayed speed is strictly closer than every competing choice
\end{definition}
\begin{proof}
  This is the declared model definition of Is Unique Closest Radial Speed Choice.
\end{proof}
% --- Archon declaration topology end ---
% --- Archon physics formalization source end ---
